\section{Introduction}

\subsection{Background and Motivation}

近年来，大语言模型（Large Language Model, LLM）的快速发展催生了一类以 LLM 为核心推理引擎的智能体（Agent）系统。这类系统通过将自然语言理解、工具调用与任务规划相结合，展现出在自动化任务处理、人机交互等领域的巨大应用潜力。在开源生态方面，LangChain\cite{chase2022langchain} 提供了面向 LLM 应用开发的模块化工具链，AutoGPT\cite{richards2023autogpt} 探索了全自主任务执行范式，MetaGPT\cite{hong2023metagpt} 则将多智能体协作引入复杂软件工程场景。与此同时，Anthropic 推出的 Computer Use\cite{anthropic2024computeruse} 方案进一步将 LLM 智能体的能力边界拓展至桌面图形界面的直接操控，标志着桌面场景智能体研究进入新的阶段。

桌面场景作为用户日常计算活动的核心承载环境，对智能体系统提出了区别于通用任务自动化的特殊需求。首先，桌面交互天然具有多模态特性，智能体需要同时处理用户文本输入与屏幕视觉信息，建立对当前桌面状态的统一感知。其次，用户的桌面操作意图多样，涵盖日常闲聊、代码编写、文件管理、系统维护等多种类型，要求智能体具备灵活的意图识别与分发能力。此外，桌面任务往往涉及多步骤的结构化工作流执行，需要在工具调用、条件判断与错误恢复之间进行精确的流程编排。

然而，现有面向桌面场景的智能体方案尚未能充分满足上述综合需求。通用 LLM 框架（如 LangChain\cite{chase2022langchain}）侧重于工具链编排，缺乏对桌面多模态上下文的深度感知；以 Computer Use\cite{anthropic2024computeruse} 为代表的方案聚焦于底层界面操控，未能构建面向用户的结构化交互流程；而 UFO\cite{zhang2024ufo}、OS-Copilot\cite{wu2024oscopilot} 等桌面智能体研究则在系统架构的分层解耦、意图分类的效率与精度平衡等方面仍有改进空间。因此，如何构建一个兼具多模态感知、高效意图识别与结构化工作流编排能力的桌面智能体架构，是当前亟需解决的关键问题。

\subsection{Problem Statement}

通过对现有桌面智能体方案的分析，可以归纳出以下四个方面的不足：

第一，缺乏结构化的分层架构。现有系统中，意图识别、角色表达与任务执行逻辑往往紧密耦合，导致各模块难以独立迭代与优化。当需要调整意图分类策略或扩展工作流节点时，通常需要对整个系统进行重构，开发与维护成本较高。

第二，意图识别策略单一。基于规则的匹配方法虽然响应速度快，但泛化能力有限，难以应对用户表达的多样性；而纯 LLM 推理方法虽具备更强的语义理解能力，却引入了显著的推理延迟。现有方案缺乏一种兼顾速度与精度的混合意图识别机制。

第三，视觉感知能力薄弱。多数桌面智能体将屏幕信息以原始图像形式直接送入 LLM 处理，未能对屏幕中的 UI 元素进行结构化识别与语义建模，无法有效利用视觉上下文来增强对话质量与任务执行的准确性。

第四，工作流编排缺乏形式化的图结构支撑。现有方案多采用线性脚本或简单的条件分支来组织任务流程，难以表达复杂的任务依赖关系与并行执行逻辑，可扩展性与容错能力不足。

\subsection{Method Overview}

针对上述问题，本文提出一种面向桌面场景的多模态分层智能体架构，采用三层解耦设计，将意图识别、角色表达与工作流执行分离为独立的功能层，以实现各模块的独立优化与灵活扩展。

路由层（Router Layer）负责用户输入的意图分类与参数抽取。该层首先基于预定义的关键词组合门控规则对输入进行快速分类，判定其所属的意图类别（如闲聊、代码任务、文件操作等）；对于需要结构化参数的任务类意图，进一步调用 LLM 进行参数抽取，将自然语言指令转化为可供下游模块消费的结构化数据。这种混合策略在保持低延迟的同时，兼顾了意图识别的准确性。

角色层（Character Layer）承担面向用户的人设驱动对话生成。该层基于预定义的角色人设（Persona）进行角色扮演式输出，生成风格一致、自然流畅的回复内容。在生成过程中，角色层同时接收来自视觉感知管线的屏幕活动上下文信息，使对话内容能够反映用户当前的桌面状态，从而提升交互的上下文感知能力与用户体验。

工作引擎层（Workflow Engine Layer）基于 LangGraph\cite{langchain2024langgraph} 框架实现图结构的工作流编排。该层将复杂任务分解为多个子图（Subgraph），通过有向图的形式显式建模任务节点之间的依赖关系、条件分支与并行执行逻辑，覆盖代码生成、文件操作、故障修复等典型桌面任务场景。

此外，本文集成了微调的 YOLOv8\cite{jocher2023yolov8} 视觉感知管线，用于实时检测屏幕中的 UI 元素并构建活动画像，将原始屏幕图像转化为结构化的视觉描述，供角色层在对话生成中使用。

\subsection{Contributions}

本文的主要贡献如下：

\begin{enumerate}
    \item 提出一种三层解耦的桌面智能体架构，将意图路由、角色扮演和工作流执行分离为独立的功能层，并给出了完整的工程实现方案，各层可独立迭代与优化。
    \item 设计并实现了一种基于关键词组合门控的快速意图分类方法，配合 LLM 结构化参数抽取，在保持低延迟的同时精准捕获用户的操作意图，兼顾了分类速度与识别精度。
    \item 基于微调的 YOLOv8\cite{jocher2023yolov8} 模型构建了桌面 UI 元素检测与活动画像映射管线，实现了对用户屏幕活动上下文的实时感知，并将视觉信息集成到角色对话生成中，增强了交互的上下文感知能力。
    \item 基于 LangGraph\cite{langchain2024langgraph} 实现了包含 5 个子图的图工作流引擎，覆盖代码生成、文件操作、故障修复等桌面任务流程，并在此基础上给出了原型系统的整体评估。
\end{enumerate}
